\section{Order and Progress}

\begin{quotex}
Yes, when the mind is missing, a soul will do. \flright{\textsc{Jean Raspail}}

\end{quotex}
Disturbed by the chaos following the French Revolution, which overthrew the order imposed by Throne and Altar, Auguste Comte sought to locate a sounder foundation for the order of society. He created a new science of man, sociology, that he hoped would support that foundation. The fundamental basis of his system is the idea of three stages of social development. One of the defects of that idea is that the three stages do not follow each other so neatly historically, but all exist at any time. The other defect is the restriction of positivism to the knowledge of all things visible. Hence, a metaphysical positivism, that also includes all tings invisible, includes Comte's insights in a larger system. Given that, the three stages can be redefined this way:

\begin{description}
\item[Mythological ]

The mythological stage is symbolic, pre-thematic, and pre-reflective, similar to what Eric Voegelin called ``compact consciousness". As pre-reflective, it exists fully at the level of being. 

\item[Ideological ]

The ideological stage regards the mythological stage as something primitive to be overcome. Thus, it attacks the foundations of the mythological society in the pursuit of ``progress", replacing the world of Being with an artificially constructed world of Thought. 

\item[Metaphysical ]

The metaphysical stage is a differentiated recapitulation of the mythological stage. Hence, it unites Being and Thought, unlike the ideological consciousness which separates them. 

\end{description}
In Comte's view, the final, or Positivistic, stage understands Progress as the evolution of Order. Evolution here is understood in its etymological sense as the ``unfolding" of what already is. In our times ``evolution" has come to mean the opposite: it means the rejection of the established order of things in order, while appealing to some ideological system that represents an allegedly higher order of being.

Now one of the appeals of ideology is its role as a status marker. It takes a certain degree of literacy and education to be initiated into an ideology. Since the mythological consciousness is pre-rational, its defenders come across as inarticulate, thus reinforcing the impressions held by ideology.

The metaphysical stage, however, is much more difficult since it requires not just a new way of thinking, but also a change in one's level of being. Anyone can adhere to an ideology: moral development or self-knowledge are not necessary to understand an ideology. Moreover, the mythological consciousness is often suspicious of the metaphysical consciousness, so it cannot usually recognize it as an ally.

We have provided several examples of the last stage. For example, they mythological and metaphysical stages roughly correspond to Guenon's distinction between the exoteric and esoteric traditions; for him, the esoteric does not reject the exoteric, but rather brings out the deeper understanding of its symbols, myths, and rites. Another example is Valentin Tomberg. Despite his deep understanding of the inner meaning of dogmas and spirituality, he regards himself as no higher than the simple believer with a symbolic, compact consciousness.

\paragraph{Roots of Order}
I have attempted here to list the roots of order of the compact society, at least of one of Solar beings. I expect the list to be modified over time. To qualify, a quality would have to be essential or innate, not something acquired over time. Hence, this is not an empirical list culled from the external behavior of several societies. As pre-rational, these beliefs are simply assumed although they are embodied in the customs, myths, institutions, etc., of the society.

Besides the natural moral law against murder, stealing, etc., there are several other properly basic qualities that are natural, without necessarily having the force of a moral law. I applied a simple test to determine the items on the list. The questions to ask are:

\begin{enumerate}
\item Is it absurd to offer an argument in support of any of the qualities? 
\item Do ideologies offer arguments opposing or attacking any of these qualities? 
\end{enumerate}
For example, for Gornahoor to post a translation of an Evola essay in support of a heterosexual nuclear family would be ludicrous. Hence, we can presume it is the default position. On the other hand, promoting some other notion of family or sexual orientation is an indication of an ideology, since it does not ``evolve" out of one of the basic positions. Here is the list of basic sources of social order, subject, of course, to comments and revision.

\begin{itemize}
\item Family as fundamental social unit 
\item Spiritual orientation and public worship 
\item Patriarchal rule 
\item Functional stratification into spiritual authority, political power, and economic activity 
\item Veneration of ancestors 
\item Preference for one's own 
\item The common good supersedes the individual 
\item Children considered a blessing rather than a burden 
\item Incest taboo 
\item Assortative mating 
\item Hypergamy 
\item Heteronormative 
\end{itemize}
\paragraph{Addendum on Dialectic and Rhetoric}
As I was getting ready to publish this, I came across this essay by \textbf{Richard Weaver}: \emph{The Cultural Role of Rhetoric}. In that essay, Mr. Weaver addresses the same point in a more sophisticated way. He opposes dialectic, or ``rational and soulless discourse", to rhetoric:

\begin{itemize}
\item \textbf{Dialectic}: abstract reasoning upon the basis of propositions. 
\item \textbf{Rhetoric}: the relation of these terms to the existential world in which facts are regarded with sympathy and are treated with that kind of historical understanding and appreciation which lie outside the dialectical process. 
\end{itemize}
Hence, in this view, dialectic is not integrated into the historical life of a community. Weaver warns us:

\begin{quotex}
Too exclusive reliance upon dialectic is a mistake of the most serious consequence because .

\end{quotex}
He writes that a society cannot live without rhetoric, wince there are some things in which the group needs to believe which cannot be demonstrated to everyone rationally. To keep questioning them (i.e., the ``roots of order") is to destroy the basis of belief in them and to weaken the cohesiveness of society. The essay concludes with an appeal, in retrospect, much too sanguine, for rhetoric and dialectic to once again join hands.

\paragraph{Enthymeme}
An enthymeme is a syllogism with one of the propositions missing. Rhetoric uses enthymeme in Aristotle's sense that

\begin{quotex}
in conversing with the multitudes you do not aim at fresh scientific instruction; you rest your arguments upon generally accepted principles and beliefs, or broadly speaking, on things received. 

\end{quotex}
Hence, the roots of order were simply assumed in public discourse with no justification needed. Weaver apparently did not foresee the extent to which those formerly generally accepted principles and beliefs have now become ``extremist". The West has not replaced rhetoric with coldly intellectual dialectic, but has instead replaced the rhetoric of being with a new, emotionally charged, neo-rhetoric. This neo-rhetoric lacks both the societal cohesiveness of rhetoric as well as scientifically derived dialectic.

Rhetoric used to serve the needs of a community with its common customs, beliefs, and historical development. Neo-rhetoric disrupts the community since it takes the view of the ``other", the outsider to the community. Neo-rhetoric has its own enthymemes which cannot be publicly disputed.



\flrightit{Posted on 2015-07-21 by Cologero }

\begin{center}* * *\end{center}

\begin{footnotesize}\begin{sffamily}



\texttt{Logres on 2015-07-22 at 07:04 said: }

``Now one of the appeals of ideology is its role as a status marker. It takes a certain degree of literacy and education to be initiated into an ideology. Since the mythological consciousness is pre-rational, its defenders come across as inarticulate, thus reinforcing the impressions held by ideology"

Politics as a form of Lust for Power: that shirt seems to fit the Progressives rather well. It's fortunate for them that their political enemies were arguing (and by and large personally believed and practiced) that power should be restricted and orderly. Also, fortunate that their enemies shared many of their revolutionary principles at some level. This made victory inevitable. It's also made them very sloppy and full of hubris. If an elite were to appear that had had an evident change of being, that would create a very resentful slum mob of progressive intelligentsia. Maybe they could Twitter their way out of it.


\hfill

\texttt{Francis on 2015-07-29 at 22:58 said: }

``Here is the list of basic sources of social order, subject, of course, to comments and revision".

Let us assume all such basics are endorsed. One must then decide to either live them in their own life alone, or make an effort to see them normalized and woven into the fabric of greater society. 

One glaring issue, to my mind, is that many of these basics become difficult to normalize in modernity as a direct result of economics. For example, one factor in the destruction of family has been the necessity of the woman to enter the work force. Dad is at ``work", mom is at ``work", and the ``family" sees one another less than they do their co-workers.

What seems an appealing option is the development of ``intentional communities", such as might be explored within the context of ``The Benedict Option". In time, what one could hope, is that these communities sprawl, and eventually ``touch" one another, thereby reinventing large territories. 

Yet, the communities still exist within the larger geography of some modern state, and ultimately remain bound by their laws; in fact, if such communities begin to gain some steam, and become a perceivable ``threat" to the existing state monopolies on power, it would be easy to predict them as targets–labeled ``cults", accused of child abuse, stockpiling weapons, evading taxes, and so on–Waco Texas serves as a reminder.

This is why, I've toyed with the idea of living Tradition within the communities, while advocating libertarianism, or ``anarcho-capitalism" without. With libertarians in outside ``public office", comes a weakening of state and enforcement power, as well as a positive deregulation of business such that larger economic wheels are greased (better facilitating trade between the Traditional intentional communities and the ``outside" world, thereby facilitating their growth and development). Thus, a bridge, or means to an end. It is often mentioned that Evola saw in the Italian Fascist state, something of a stepping stone toward a world of Tradition–it was a model that suited his times, and existed. Today, Fascism of that Italian model does not exist, nor will return again–so, little good in attempting to ``restore" that, or advance a ``Mussolinism". Nor shall we see, some sudden return to monarchies, built from the tattered remnants of living hereditary heirs. 

Instead, a new, fresh, and practical model must be considered as a ``bridge"–that is, if one is inclined to a path of ``action". And, that bridge, I reiterate, is outward advancement of candidates representing an anarcho-capitalism/libertarianism, while internally building self-sustaining intentional communities having Tradition at their core, governed along aristocratic principles (unlike leftist ``communes"), out of which, an ``elite" will begin to take shape.


\end{sffamily}\end{footnotesize}
